\section{Conclusion}\label{conclusion}

We presented a ConvNeXt-Tiny U-shaped segmentation framework that uses a locally
implemented bottleneck FAFEM and couples its descriptor to the MSCB structure through
RFG-MSCB at the Stage-3 encoder--decoder fusion. The design
extends bottleneck frequency enhancement into decoder refinement while preserving the
ordinary multi-scale pathway as a reference, thereby addressing the motivation for
sample-dependent receptive-field selection without replacing the established decoder
path.

Across three seeds, the model achieved the highest displayed Kvasir-SEG mDice and
mIoU, with strong performance also observed on CVC-ClinicDB. The parameter-efficient
modification is localized to Stage 3, adds 2.4\% parameters relative to FAFEM for a
30.06 M-parameter architecture, and retains depthwise-separable convolution in the
remaining decoder.

Future work should examine placement across multiple seeds, separate residualization
from frequency conditioning, and evaluate domain-adaptation strategies for the
challenging external datasets. Overall, RFG-MSCB connects bottleneck frequency
information to sample-dependent multi-scale decoding through bounded residual coupling
that preserves the original MSCB pathway.
